
\documentclass{article}
\usepackage{colortbl}
\usepackage{makecell}
\usepackage{multirow}
\usepackage{supertabular}

\begin{document}

\newcounter{utterance}

\centering \large Interaction Transcript for game `bbh', experiment `dyck\_languages', episode 87 with reflexion{-}Qwen3.5{-}27B.
\vspace{24pt}

{ \footnotesize  \setcounter{utterance}{1}
\setlength{\tabcolsep}{0pt}
\begin{supertabular}{c@{$\;$}|p{.15\linewidth}@{}p{.15\linewidth}p{.15\linewidth}p{.15\linewidth}p{.15\linewidth}p{.15\linewidth}}
    \# & \multicolumn{2}{c}{Player} && \multicolumn{2}{c}{Game Master} \\
    \hline

    \theutterance \stepcounter{utterance}  
    & & & \multicolumn{4}{p{0.6\linewidth}}{
        \cellcolor[rgb]{0.9,0.9,0.9}{
            \makecell[{{p{\linewidth}}}]{
                \texttt{\tiny{[P1$\langle$GM]}}
                \texttt{You have attempted this same task before and did not succeed. These are the plans you wrote after those attempts.} \\
\texttt{Trial \#1: New plan: My previous attempt failed because I ignored the specific types of brackets in the input and simply outputted a long string of closing square brackets `{]}`. This is incorrect because a Dyck word requires matching opening and closing pairs of the same type (e.g., `{[}` must be closed by `{]}`, `\{` by `\}`, `<` by `>`, and `(` by `)`). To succeed, I must parse the input string from left to right, maintaining a stack of open brackets. For every opening bracket encountered, I will push it onto the stack. Once the input string is fully processed, I will pop the brackets from the stack one by one and append their corresponding closing brackets to the output in reverse order of their appearance. This ensures that the most recently opened bracket is closed first, resulting in a valid, balanced sequence.} \\
\\ 
\texttt{Trial \#2: New plan: My previous attempt failed because I stopped the stack analysis prematurely at step 19, leaving a significant portion of the input string unprocessed. This resulted in an incomplete stack state and an incorrect output. To succeed, I must process the **entire** input string character by character without interruption. I will maintain a stack, pushing every opening bracket (`{[}`, `\{`, `<`, `(`) and popping the matching closing bracket (`{]}`, `\}`, `>`, `)`) when encountered. Only after the entire input string has been fully traversed will I generate the output by popping the remaining items from the stack in reverse order and appending their corresponding closing brackets. This ensures all nested structures are correctly balanced.} \\
\\ 
\texttt{Trial \#3: New plan: My previous attempt failed because I abandoned the stack analysis midway through the input string (stopping at step 17), leaving a large portion of the sequence unprocessed. This resulted in an incomplete stack state and an incorrect final output. To succeed, I must rigorously process the **entire** input string character by character without stopping early. I will maintain a stack, pushing every opening bracket (`{[}`, `\{`, `<`, `(`) and popping the matching closing bracket (`{]}`, `\}`, `>`, `)`) when encountered. Only after the full string is traversed will I generate the output by popping the remaining items from the stack in reverse order (LIFO) and appending their corresponding closing brackets. I will ensure the final output is a single string of closing brackets with no explanatory text or intermediate steps included.} \\
\\ 
\\ 
\texttt{Correctly close a Dyck{-}n word.} \\
\\ 
\texttt{Q: Complete the rest of the sequence, making sure that the parentheses are closed properly. Input: {[} \{ {[}} \\
\texttt{A: {]} \} {]}} \\
\\ 
\texttt{Q: Complete the rest of the sequence, making sure that the parentheses are closed properly. Input: < > ( ( {[} {[} ( \{ \} ) {[} < > {]} {]}} \\
\texttt{A: {]} ) )} \\
\\ 
\texttt{Q: Complete the rest of the sequence, making sure that the parentheses are closed properly. Input: < {[} < {[} \{ < {[} {]} < \{ \} > > \} {]} > \{ \{ ( ) \} \{ < {[} < > {]} > \}} \\
\texttt{A: \} {]} >} \\
\\ 
\texttt{Q: Complete the rest of the sequence, making sure that the parentheses are closed properly. Input: {[} {[} < < \{ {[} {]} ( {[} {[} ( {[} {[} \{ {[} {]} \} {]} {]} ) ( ) {]} {]} ) ( ) {[} {]} < {[} ( ( ( ) ) ) {]} < < < > {[} ( ) {]} > {[} {]} > > {[} < > {]} \} > >} \\
\texttt{A:} \\
            }
        }
    }
     \\ \\

    \theutterance \stepcounter{utterance}  
    & \multicolumn{4}{p{0.6\linewidth}}{
        \cellcolor[rgb]{0.561,0.737,0.561}{
            \makecell[{{p{\linewidth}}}]{
                \texttt{\tiny{[P1$\rangle$GM]}}
                \texttt{{]} \} > > {]}} \\
            }
        }
    }
    & & \\ \\

    \theutterance \stepcounter{utterance}  
    & & & \multicolumn{2}{p{0.3\linewidth}}{
        \cellcolor[rgb]{0.9,0.9,0.9}{
            \makecell[{{p{\linewidth}}}]{
                \texttt{\tiny{[GM$|$GM]}}
                \texttt{Target: {]} {]}} \\
            }
        }
    }
    & & \\ \\

    \theutterance \stepcounter{utterance}  
    & & & \multicolumn{2}{p{0.3\linewidth}}{
        \cellcolor[rgb]{0.9,0.9,0.9}{
            \makecell[{{p{\linewidth}}}]{
                \texttt{\tiny{[GM$|$GM]}}
                \texttt{game\_result = LOSE} \\
            }
        }
    }
    & & \\ \\

    \theutterance \stepcounter{utterance}  
    & & \multicolumn{4}{p{0.6\linewidth}}{
        \cellcolor[rgb]{0.95,0.95,0.95}{
            \makecell[{{p{\linewidth}}}]{
                \texttt{\tiny{[P1$\leftrightarrow$P]}}
                \texttt{Post{-}episode Reflexion (not sent to GM):} \\
\texttt{New plan: My previous attempts failed because I did not rigorously process the entire input string character by character, leading to an incomplete stack state and an incorrect sequence of closing brackets. Specifically, I likely stopped analyzing the input prematurely or miscounted the nesting depth in the complex final example. To succeed, I will simulate a stack data structure mentally for the **entire** input string: I will push every opening bracket (`{[}`, `\{`, `<`, `(`) onto the stack and pop the top element only when a matching closing bracket is encountered. Once the input string is fully exhausted, I will generate the output by popping the remaining elements from the stack in Last{-}In{-}First{-}Out (LIFO) order and appending their corresponding closing brackets (`{]}`, `\}`, `>`, `)`) to form the completion. I will ensure no characters are skipped and the final output contains only the necessary closing brackets.} \\
            }
        }
    }
    & \\ \\

\end{supertabular}
}

\end{document}
